
% ============================================================================
\clearpage
\section*{v4.10.0-r12 Source-Only: Status Notes on Inherited Koide and Generation Statements}
\addcontentsline{toc}{section}{v4.10.0-r12 Source-Only: Status Notes on Inherited Koide and Generation Statements}
\label{sec:v4100-r12-status-notes}
% ============================================================================

\paragraph{Purpose.}
Several inherited passages state more about the Koide block and the generations than the later layers support.  They were
noted during the OP3 work after v5.0.0 (12 and 13 November 2026, notes under \texttt{notes/\allowbreak research/\allowbreak op3\_*}) and
are recorded here as status notes, as the append-only rule requires.  The layer also registers the two neutrino tests of
that work (H99 evaluated, H100 pre-registered) and the explorations behind them.  No new test is evaluated here.  The notes
are \texttt{notes/\allowbreak research/\allowbreak v4\_10\_0\_r12/}.  Earlier layers are not rewritten.

\noindent\textbf{Source-only status.}
This v4.10.0-r12 layer is source-only.  It does not build a PDF, does not assign a new DOI,
and does not authorize a public release.  The v5.0.0 Version DOI
\href{https://doi.org/10.5281/zenodo.24152615}{10.5281/zenodo.24152615} and the Concept DOI
\href{https://doi.org/10.5281/zenodo.19210728}{10.5281/zenodo.19210728} are unchanged by it.

\subsection*{v4.10.0-r12 status notes}
\label{subsec:v4100-r12-status}

\paragraph{Theoremlet v4.10.0-r12.1 (historical OP~3 is not closed).}
The open-question genealogy (part 17, ``Current reading'' and ``Historical OP synchronization'') calls the Koide/triolic block
``structurally closed at the present level'' and says that historical OP~3 should no longer be listed as a primary open problem.
This is superseded.  The Koide relation is an empirical fact whose derivation is open (status of
Theorem~\ref{thm:koide-triolic}); the phase \(\delta=2/9\) holds only for charged-lepton pole masses and fails for quarks and
for running lepton masses (v4.8.0-r11, H1); the overlap model of the v4.11 section is conditional; the v5.0.0 string question
is answered negatively only for the Maxwell class and stays open in general.  OP~3 remains open.

\paragraph{Theoremlet v4.10.0-r12.2 (the remark on the factor \(2/3\)).}
The remark ``The \(2/3\) factor is everywhere'' lists four occurrences of \(2/3\) and reads them as one statement (two of the
three triolic sectors collapse to the photon ground state).  No common derivation exists.  The measured dark-energy fraction
is about \(0.685\), not \(2/3\); the spring coefficient and the stable sector fraction are separate constructions; for the Koide
ratio the Gram reading of the post-v5 notes gives \(K=(1+2\rho^2)/3\), so \(K=2/3\) is a mean squared overlap
\(\rho^2=1/2\), an assumption-bearing readout and not a sector count.  The remark is a list of coincidences and is not to be
read as a result.

\paragraph{Theoremlet v4.10.0-r12.3 (remark on the mass discrepancy near the electron).}
Remark~\ref{rem:koide-mass-discrepancy} attributes the residual mass errors to the approximation of \(\beta_{\rm sym}\) in the
phase formula \(\theta_K=\pi-\sqrt2\arcsin\beta_{\rm sym}\).  That formula is excluded at about \(300\)--\(400\sigma\)
(v4.8.0-r09) and archived (status of Theorem~\ref{thm:koide-triolic}).  The attribution is therefore void.  The amplification
near the electron (a small change of the phase moves \(m_e\) strongly) remains a correct property of the Brannen map.

\paragraph{Theoremlet v4.10.0-r12.4 (generation labels).}
\begin{enumerate}[leftmargin=2em]
\item[(a)] The left-invariant Gram labels \(\mu_i\) of v4.9.0-r33 are not invariant under the electromagnetic gauge
transformation: right multiplication in the quaternion factor is the weak isospin, \(Q_{\rm em}\) contains \(T_3\), a
\(Q_{\rm em}\) rotation mixes \(i\) and \(j\), so \(\mu_1\) and \(\mu_2\) change while \(\mu_3\) does not.  Only the eigenvalues of
the Gram matrix are charge invariant (note \texttt{op3\_\allowbreak generation\_\allowbreak carrier\_\allowbreak review\_\allowbreak 2026\_\allowbreak 11\_\allowbreak 13.md};
the assumption is the same quaternion factor as in v4.9.0-r40 to r43).  The refutation of H18 is unchanged.
\item[(b)] Theorem~\ref{thm:three-generations} keeps its v4.9.0 status (the triplet basis is colour).  In addition, none of the
carriers examined after v5.0.0 gives the number three and the masses at once (colour triplet, quaternion directions, Hopf
carrier, place-index \(U(3)\), sedenion plane); a second three-fold structure outside \(\OO\) is still required.
\item[(c)] Exact, found in the explorations: three spin-\(\tfrac12\) states along three orthogonal axes have pairwise overlap
\(\tfrac12\), rank \(2\) and Bargmann phase \(\pi/4\); this is precisely the boundary triple of the Koide arc
(\(\rho^2=\tfrac12\), \(\theta=\pi/12\)).  The cyclic right action of \(1+i+j+k\) in r33.2 is, up to normalisation, the
Hurwitz unit that generates this triple.  No colour- and charge-neutral native carrier was found (isospin breaks
\(Q_{\rm em}\), octonionic directions carry colour, the neutral Dixon part has signature \((2,1)\), the sedenion cycle \(\psi\)
moves the charge axis).
\end{enumerate}

\paragraph{Theoremlet v4.10.0-r12.5 (neutrino tests of the post-v5 work).}
\begin{enumerate}[leftmargin=2em]
\item[(a)] H99 (pre-registration \texttt{5b23c5a}, not blind): equal overlap \(\rho^2=\tfrac12\) in both colour-neutral sectors and
\(\theta_\nu=\theta_e+\pi/12\) give \(R=\Delta m^2_{21}/\Delta m^2_{31}=0.030839\).  Against NuFIT~6.0 \(+1.31\sigma\), against
NuFIT~6.1 (primary table, IC24 with SK, normal ordering) \(+1.80\sigma\): compatible by the criterion.
\item[(b)] H100 (pre-registration \texttt{9c77ddd}): three values of \(R\) are fixed for the next NuFIT release: \(0.030839\)
(H99), \(0.030087\) (all ingredients at one short-distance scale), \(0.030375\) (Sumino-type family correction, \(c=1\),
PMNS-averaged).  Pending.
\item[(c)] Explorations (not tests): the Brannen phase is one third of the Bargmann phase of the three places; H99 is a flux
\(\pi/4\) through the place triangle; the native phases \(\omega\), \(i\), \(-1\) generate only multiples of \(\pi/6\), so
\(\pi/4\) needs a further source; a rule ``phase \((Y/2\text{ of the connector})\cdot\pi/6\)'' reproduces \(\pi/12\) for
Majorana neutrinos (\(HH\)) and \(\pi/6\) for Dirac neutrinos, a posteriori and without mechanism.
\end{enumerate}

\subsection*{v4.10.0-r12 status in the three frames}
\label{subsec:v4100-r12-frames}

\begin{itemize}[leftmargin=2em]
\item M and T: the status notes concern readouts in the matter frame; the Gram and Bargmann statements are frame free in the
places; nothing is transported to T.
\item P: not evaluated.
\end{itemize}

\subsection*{v4.10.0-r12 register entries}
\label{subsec:v4100-r12-register}

\begin{center}
\small
\begin{tabular}{@{}>{\raggedright\arraybackslash}p{0.34\textwidth}>{\raggedright\arraybackslash}p{0.58\textwidth}@{}}
\toprule
statement & status after v4.10.0-r12\\
\midrule
historical OP~3 ``structurally closed'' (part 17) & superseded; OP~3 open (r12.1)\\
remark on the factor \(2/3\) & list of coincidences, not a result (r12.2)\\
Remark~\ref{rem:koide-mass-discrepancy} & attribution to \(\beta_{\rm sym}\) void (r12.3)\\
r33 Gram labels \(\mu_i\) & not \(Q_{\rm em}\) invariant (r12.4(a))\\
Koide arc boundary & octant triple of spin-\(\tfrac12\) states, exact (r12.4(c))\\
H99 / H100 & compatible \(+1.80\sigma\) / pending (r12.5)\\
\bottomrule
\end{tabular}
\end{center}

\subsection*{v4.10.0-r12 ledger and guard update}
\label{subsec:v4100-r12-ledger-guard}

\[
R_{\rm chk}^{v4100\text{-}r12}=(1,1,0)
\]
for (status notes recorded; formulas checked in the cited notes; native carrier and derivation of \(K\) open).

The layer spends no guard:
\[
R_{\rm glob}^{(v4100\text{-}r12,+)}=(0,1,1,1,1),
\qquad
\Gamma_{\rm guard}^{v4100\text{-}r12}=(1,1,1,1,1),
\qquad
\beta_4=1/256\ \text{unspent}.
\]
It does not imply exact finite \(C5\) zero closure, exact all-readout terminality,
\(C6/T2\), a global photon-clock/RH theorem, or global all-readout closure.
